\documentclass[a4paper,12pt]{article}
\usepackage[utf8]{inputenc}
\usepackage[T2A]{fontenc}
\usepackage[russian]{babel}
\usepackage{amsmath, amssymb, amsfonts}
\usepackage{geometry}
\usepackage{xcolor}
\usepackage{hyperref}

\geometry{left=2cm, right=2cm, top=2cm, bottom=2cm}

% ============================================================
% ЦВЕТОВОЕ ВЫДЕЛЕНИЕ МОДИФИКАЦИЙ
% Измените цвет здесь для всех модификаций сразу
% ============================================================
\definecolor{modcolor}{rgb}{0.0, 0.55, 0.0}  % зелёный (измените по желанию)
% \definecolor{modcolor}{rgb}{0.0, 0.0, 0.7}  % синий (альтернатива)
\definecolor{redcolor}{rgb}{0.7, 0.0, 0.0}  % красный (альтернатива)

\newcommand{\redtext}[1]{\textcolor{redcolor}{#1}}
\newcommand{\redmath}[1]{\textcolor{redcolor}{#1}}
\newcommand{\redblock}[1]{\colorbox{redcolor!10}{\parbox{\dimexpr\linewidth-2\fboxsep}{#1}}}

\newcommand{\modtext}[1]{\textcolor{modcolor}{#1}}
\newcommand{\modmath}[1]{\textcolor{modcolor}{#1}}
\newcommand{\modblock}[1]{\colorbox{modcolor!10}{\parbox{\dimexpr\linewidth-2\fboxsep}{#1}}}

\title{Пондеромоторные силы в диэлектриках и магнетиках:\\
Модифицированный вывод И.Е.~Тамма для полной системы уравнений Максвелла\\
с учётом динамики, магнитных токов и разрывов на границах сред}
\author{И.Е.~Тамм (\S 17, 18, 32--34, 84, 105)\\
с модификациями для модели MenDrive (А.Ю.~Дроздов)}
\date{}

\begin{document}

\maketitle

\begin{abstract}
Настоящий документ представляет собой подробное воспроизведение вывода пондеромоторных 
сил и тензора натяжений электромагнитного поля из книги И.Е.~Тамма «Основы теории 
электричества» с критическими модификациями, необходимыми для применения к модели 
волнового двигателя MenDrive. 

\modtext{
    \textbf{Модификации вносятся только в тех местах, где оригинальный вывод Тамма 
    содержит предположения, не выполняющиеся для MenDrive:}
    \begin{enumerate}
        \item Динамические поля (формула 18.9 вместо 17.5)
        \item Магнитные заряды и токи ($\rho_m \neq 0$, $\vec{j}_m \neq 0$)
        \item Разрывы $\varepsilon$, $\mu$ на границах сред (теорема Гаусса применяется 
        отдельно к каждой области)
        \item Поверхностные интегралы на внутренних границах не обращаются в ноль
    \end{enumerate}
}

Все модификации выделены \textcolor{modcolor}{зелёным цветом}. Оригинальный текст 
Тамма приведён чёрным цветом.
\end{abstract}

\tableofcontents

\newpage

% ============================================================
% §17. ПОНДЕРОМОТОРНЫЕ СИЛЫ (стр. 88–91)
% ============================================================

\section*{\S\,17. Пондеромоторные силы.}

\subsection*{1. Сила на электрический диполь}

\noindent
ников расположены на их поверхностях, так что $\rho = 0$, и что 
потенциал каждого проводника постоянен на всём его протяжении, 
получим из

\begin{equation}
F = fS = \frac{2\pi e^2}{S} = 2\pi\sigma^2\cdot S,
\end{equation}
где $\sigma = \dfrac{e}{S}$ — плотность заряда на поверхности пластин, 
что вполне согласуется с (17.2).

К тому же результату можно прийти и непосредственно из (18.3), 
полагая в этой формуле $q = d$ и $\Theta = F$.

\textbf{3.} В предыдущем мы неявным образом предполагали, что при 
смещении пластин конденсатора \textit{заряд} их $e$ остаётся 
постоянным. Мы получили бы совершенно иное

\modtext{
    \textbf{МОДИФИКАЦИЯ 17.1:} Для динамических полей необходимо использовать 
    полную формулу силы на диполь, которая включает члены с производными по времени. 
    В оригинальном выводе Тамма используется статическая формула (17.5), но для 
    MenDrive с гармоническими полями $e^{-i\omega t}$ это недопустимо.
}

\modblock{
\modtext{
    Для \textbf{динамических полей} сила на электрический диполь определяется 
    формулой (18.9) вместо (17.5):
    \begin{equation}
    \modmath{
        \vec{F} = (\vec{p}\cdot\nabla)\vec{E} + 
        \left[\vec{p} \times \text{rot}\,\vec{E}\right] +
        \frac{1}{c}\left[\frac{\partial\vec{p}}{\partial t} \times \vec{B}\right]
    }
    \tag{17.5$^*$}
    \end{equation}
    Последний член отсутствует в статической формуле (17.5) Тамма, но для 
    гармонических полей $\vec{p}(t) = \vec{p}_0 e^{-i\omega t}$ он даёт вклад:
    \begin{equation}
    \modmath{
        \frac{1}{c}\left[\frac{\partial\vec{p}}{\partial t} \times \vec{B}\right] = 
        -\frac{i\omega}{c}\left[\vec{p} \times \vec{B}\right]
    }
    \end{equation}
    Этот член ответственен за \textbf{Абрагам-подобную силу} $\vec{F}_P$ в расчёте 
    MenDrive.
}
}

\subsection*{2. Плотность сил в диэлектрике}

\noindent
Плотность сил, приложенных к диэлектрику

\begin{equation}
\vec{f} = \left(\vec{P}\nabla\right)\cdot\vec{E} + \left[\vec{P} \times \text{rot}\,\vec{E}\right]
\tag{17.6}
\end{equation}

\modtext{
    \textbf{МОДИФИКАЦИЯ 17.2:} Для полной системы с магнитными токами необходимо 
    добавить магнитную часть и динамические члены:
}

\modblock{
\modtext{
    \begin{equation}
    \modmath{
    \vec{f} = \left(\vec{P}\nabla\right)\cdot\vec{E} + 
    \left[\vec{P} \times \text{rot}\,\vec{E}\right] +
    \frac{1}{c}\left[\frac{\partial\vec{P}}{\partial t} \times \vec{B}\right] +
    \left(\vec{M}\nabla\right)\cdot\vec{H} +
    \left[\vec{M} \times \text{rot}\,\vec{H}\right] -
    \frac{1}{c}\left[\frac{\partial\vec{M}}{\partial t} \times \vec{D}\right]
    }
    \tag{17.6$^*$}
    \end{equation}
}
}

\noindent
где поляризация диэлектрика равна

\begin{equation}
\vec{P} = \frac{\varepsilon - 1}{4\pi}\vec{E}
\tag{17.7}
\end{equation}

поэтому

\begin{equation}
\vec{f} = \left(\frac{\varepsilon - 1}{4\pi}\vec{E}\nabla\right)\cdot\vec{E} + 
\left[\frac{\varepsilon - 1}{4\pi}\vec{E} \times \text{rot}\,\vec{E}\right]
\tag{17.8}
\end{equation}

\modtext{
    \textbf{МОДИФИКАЦИЯ 17.3:} В модели MenDrive тензор диэлектрической проницаемости 
    \textit{диагонален} но не скалярен: $\varepsilon_{xx} \neq \varepsilon_{yy} \neq \varepsilon_{zz}$. 
    Кроме того, в проводниках $\varepsilon$ комплексна из-за проводимости $\sigma_e$.
}

\modblock{
\modtext{
    \begin{equation}
    \modmath{
    \vec{f} = \left(\frac{\overleftrightarrow{\varepsilon} - \hat{I}}{4\pi}\vec{E}\nabla\right)\cdot\vec{E} + 
    \left[\frac{\overleftrightarrow{\varepsilon} - \hat{I}}{4\pi}\vec{E} \times \text{rot}\,\vec{E}\right] +
    \frac{1}{c}\left[\frac{\partial}{\partial t}\left(\frac{\overleftrightarrow{\varepsilon} - \hat{I}}{4\pi}\vec{E}\right) \times \vec{B}\right]
    }
    \tag{17.8$^*$}
    \end{equation}
}
}

\subsection*{3. Компонента силы вдоль оси $x$}

\noindent
Рассмотрим слагающую плотности сил $\vec{f}$ по какому-нибудь направлению, 
например по оси $x$:

\begin{equation}
f_x = P_x\frac{\partial E_x}{\partial x} + 
P_y\frac{\partial E_y}{\partial x} + 
P_z\frac{\partial E_z}{\partial x} + 
\left[\vec{P} \times \text{rot}\,\vec{E}\right]_x
\tag{17.9}
\end{equation}

\modtext{
    \textbf{МОДИФИКАЦИЯ 17.4:} Подставляем $\vec{P} = \frac{\varepsilon - 1}{4\pi}\vec{E}$ 
    и учитываем тензорный характер $\varepsilon$:
}

\modblock{
\modtext{
    \begin{equation}
    \modmath{
        \begin{aligned}
        f_x &= \left(\frac{\varepsilon_{xx} - 1}{4\pi}E_x\right)\frac{\partial E_x}{\partial x} + 
        \left(\frac{\varepsilon_{yy} - 1}{4\pi}E_y\right)\frac{\partial E_y}{\partial x} + 
        \left(\frac{\varepsilon_{zz} - 1}{4\pi}E_z\right)\frac{\partial E_z}{\partial x} + \\
        &\quad + \left[\left(\frac{\overleftrightarrow{\varepsilon} - 1}{4\pi}\vec{E}\right) \times \text{rot}\,\vec{E}\right]_x +
        \frac{1}{c}\left[\frac{\partial}{\partial t}\left(\frac{\overleftrightarrow{\varepsilon} - 1}{4\pi}\vec{E}\right) \times \vec{B}\right]_x
        \end{aligned}
    }
    \tag{17.9$^*$}
    \end{equation}
}
}

\subsection*{4. Интегрирование по глубине материала}

\modtext{
    \textbf{МОДИФИКАЦИЯ 17.5:} Тамм не выполняет явно это интегрирование, но для 
    MenDrive критически важно проинтегрировать по глубине проникновения поля в 
    проводник (скин-слой). При условии полного затухания поля на бесконечности:
}

\modblock{
\modtext{
    \begin{equation}
    \modmath{
        \begin{aligned}
        \int_{a}^{\infty} f_x\,dx &= 
        \left(\frac{\varepsilon_{xx} - 1}{4\pi}E_x\right)\frac{E_x}{2}\Bigg|_{a}^{\infty} +
        \left(\frac{\varepsilon_{yy} - 1}{4\pi}E_y\right)\frac{E_y}{2}\Bigg|_{a}^{\infty} +
        \left(\frac{\varepsilon_{zz} - 1}{4\pi}E_z\right)\frac{E_z}{2}\Bigg|_{a}^{\infty} + \\
        &\quad + \int_{a}^{\infty}\left[\left(\frac{\overleftrightarrow{\varepsilon} - 1}{4\pi}\vec{E}\right) \times \text{rot}\,\vec{E}\right]_x dx +
        \int_{a}^{\infty}\frac{1}{c}\left[\frac{\partial}{\partial t}\left(\frac{\overleftrightarrow{\varepsilon} - 1}{4\pi}\vec{E}\right) \times \vec{B}\right]_x dx
        \end{aligned}
    }
    \tag{17.10$^*$}
    \end{equation}
    При $E|_{x\to\infty} = 0$ получаем:
    \begin{equation}
    \modmath{
        \begin{aligned}
        \int_{a}^{\infty} f_x\,dx =
        -\left(\frac{\varepsilon_{xx} - 1}{4\pi}E_x(x=a)\right)\frac{E_x(x=a)}{2} -\\
        \left(\frac{\varepsilon_{yy} - 1}{4\pi}E_y(x=a)\right)\frac{E_y(x=a)}{2} -
        \left(\frac{\varepsilon_{zz} - 1}{4\pi}E_z(x=a)\right)\frac{E_z(x=a)}{2} + \ldots
        \end{aligned}
    }
    \tag{17.11$^*$}
    \end{equation}
}
}

\noindent
при условии полного угасания поля в скин-слое первые три слагаемые приводят 
к формулам давления создаваемым электрическим полем из курса электростатики Смайта

\begin{equation}
\begin{aligned}
\int f_x\,dx =
\left(\frac{1 - \varepsilon_{xx}}{4\pi}E_x(x=a)\right)\frac{E_x(x=a)}{2} + \\
\left(\frac{1 - \varepsilon_{yy}}{4\pi}E_y(x=a)\right)\frac{E_y(x=a)}{2} +
\left(\frac{1 - \varepsilon_{zz}}{4\pi}E_z(x=a)\right)\frac{E_z(x=a)}{2} + \ldots
\end{aligned}
\tag{17.12}
\end{equation}

% ============================================================
% §18. ДИНАМИЧЕСКИЕ ПОПРАВКИ (стр. 92–95)
% ============================================================

\section*{\S\,18. Динамические поправки к пондеромоторным силам.}

\modtext{
    \textbf{ПРИМЕЧАНИЕ:} Этот параграф отсутствует в оригинальном выводе Тамма, 
    но критически важен для MenDrive. В §105 Тамм упоминает член с импульсом 
    поля $\vec{g}$, но не включает его явно в вывод тензора.
}

\subsection*{1. Полная плотность силы для динамических полей}


% \modblock{
% \modtext{

    Для гармонических полей $\vec{E}(t) = \vec{E}_0 e^{-i\omega t}$, 
    $\vec{H}(t) = \vec{H}_0 e^{-i\omega t}$ полная плотность пондеромоторной 
    силы включает \textbf{динамические члены}:

    \begin{equation}
    \modmath{
        \vec{f}_{\text{full}} = \vec{f}_{\text{static}} + \vec{f}_{\text{dynamic}}
    }
    \tag{18.1}
    \end{equation}

    где статическая часть:
    \begin{equation}
    \modmath{
        \vec{f}_{\text{static}} = \rho_e\vec{E} + \frac{1}{c}[\vec{j}_e\times\vec{B}] + 
        \rho_m\vec{H} + \frac{1}{c}[\vec{j}_m\times\vec{D}] + 
        (\vec{P}\cdot\nabla)\vec{E} + (\vec{M}\cdot\nabla)\vec{H}
    }
    \tag{18.2}
    \end{equation}

    и динамическая часть (отсутствует у Тамма в явном виде):
    \begin{equation}
    \modmath{
        \vec{f}_{\text{dynamic}} = \frac{1}{c}\left[\frac{\partial\vec{P}}{\partial t}\times\vec{B}\right] - 
        \frac{1}{c}\left[\frac{\partial\vec{M}}{\partial t}\times\vec{D}\right]
    }
    \tag{18.3}
    \end{equation}

    \subsection*{2. Усреднение по периоду}

    Для гармонических полей среднее за период значение:
    \begin{equation}
    \modmath{
        \langle\vec{f}\rangle = \frac{1}{T}\int_0^T \vec{f}(t)\,dt = 
        \frac{1}{2}\text{Re}\left[\vec{f}_0\right]
    }
    \tag{18.4}
    \end{equation}

    где $\vec{f}_0$ — комплексная амплитуда. Это даёт множитель $1/2$ в формулах 
    для компонент силы в коде MenDrive:
    \begin{equation}
    \modmath{
        \text{additional 1/2 arises as result of integration on period}
    }
    \tag{18.5}
    \end{equation}

% }
% }

% ============================================================
% §32. ПОНДЕРОМОТОРНЫЕ СИЛЫ В ДИЭЛЕКТРИКАХ (стр. 148–168)
% ============================================================

\section*{\S\,32. Пондеромоторные силы в диэлектриках.}

\subsection*{1. Разложение плотности сил}

\noindent
С этой целью удобно разложить общее выражение (32.10) для объёмной 
плотности этих сил на два слагаемых:

\begin{equation}
\vec{f} = \vec{f}' + \vec{f}'',
\quad
\vec{f}' = \rho\vec{E} - \frac{1}{8\pi}E^2\operatorname{grad}\varepsilon,
\quad
\vec{f}'' = \frac{1}{8\pi}\operatorname{grad}\!\left(
E^2\frac{\partial\varepsilon}{\partial\tau}\tau
\right).
\tag{32.1}
\end{equation}

\modtext{
\textbf{МОДИФИКАЦИЯ 32.1:} В оригинальном выводе Тамма предполагается 
$\rho = 0$ (нет свободных зарядов). Но в проводниках MenDrive 
$\rho_e \neq 0$ из-за $\text{div}\,\vec{D} \neq 0$ в динамике. Также 
отсутствует магнитная часть с $\rho_m$, $\vec{j}_m$.
}

% \modblock{
% \modtext{
\begin{equation}
\modmath{
\vec{f} = \vec{f}'_e + \vec{f}''_e + \vec{f}'_m + \vec{f}''_m + \vec{f}_{\text{dynamic}}
}
\tag{32.1$^*$}
\end{equation}
где:
\begin{align}
% \modmath{
    \vec{f}'_e &= \rho_e\vec{E} - \frac{1}{8\pi}E^2\operatorname{grad}\varepsilon \\
    \vec{f}''_e &= \frac{1}{8\pi}\operatorname{grad}\!\left(
    E^2\frac{\partial\varepsilon}{\partial\tau}\tau
    \right) \\
    \vec{f}'_m &= \rho_m\vec{H} - \frac{1}{8\pi}H^2\operatorname{grad}\mu \\
    \vec{f}''_m &= \frac{1}{8\pi}\operatorname{grad}\!\left(
    H^2\frac{\partial\mu}{\partial\tau}\tau
    \right) \\
    \vec{f}_{\text{dynamic}} &= \frac{1}{c}\left[\frac{\partial\vec{P}}{\partial t}\times\vec{B}\right] - 
    \frac{1}{c}\left[\frac{\partial\vec{M}}{\partial t}\times\vec{D}\right]
% }
\end{align}
% }
% }

\subsection*{2. Выражение через поля}

\noindent
Наша задача будет, очевидно, разрешена, если мы найдём такой тензор 
$\mathbf{t}$, чтобы по подстановке его компонент в правую часть 
уравнений (33.7) левые части этих уравнений совпали с определяемыми 
формулой (34.1) слагаемыми плотности объёмных сил $\vec{f}'$.

\textbf{2.} Выразим в (34.1) $\rho$ через $\operatorname{div}\vec{D}$ 
с помощью (22.2) и рассмотрим слагающую плотности сил $\vec{f}'$ 
по какому-нибудь направлению, например по оси $x$:

\begin{equation}
f_x = \frac{1}{4\pi}E_x\operatorname{div}\vec{D} - 
\frac{1}{8\pi}E^2\frac{\partial\varepsilon}{\partial x}
\tag{32.2}
\end{equation}

\modtext{
\textbf{МОДИФИКАЦИЯ 32.2:} Тамм использует $\text{div}\,\vec{D} = 4\pi\rho$, 
но в динамике для проводников $\text{div}\,\vec{D} \neq 0$ даже при 
отсутствии свободных зарядов из-за конечной проводимости $\sigma_e \neq 0$.
}

\modblock{
\modtext{
Из уравнений Максвелла с магнитными токами:
\begin{equation}
\modmath{
\operatorname{div}\vec{D} = 4\pi\rho_e, \quad
\operatorname{div}\vec{B} = 4\pi\rho_m
}
\tag{32.2$^*$}
\end{equation}
Для гармонических полей в проводнике:
\begin{equation}
\modmath{
\rho_e = \frac{1}{4\pi}\operatorname{div}\vec{D} = 
\frac{\varepsilon}{4\pi}\operatorname{div}\vec{E} + 
\frac{1}{4\pi}\vec{E}\cdot\operatorname{grad}\varepsilon
}
\tag{32.3$^*$}
\end{equation}
В модели MenDrive $\operatorname{grad}\varepsilon \neq 0$ на границах 
сред, что даёт поверхностные члены.
}
}

\subsection*{3. Преобразование к дивергенции тензора}

\noindent
Преобразуем первое слагаемое:

\begin{equation}
E_x\operatorname{div}\vec{D} = 
\frac{\partial}{\partial x}(E_x D_x) + 
\frac{\partial}{\partial y}(E_x D_y) + 
\frac{\partial}{\partial z}(E_x D_z) - 
\vec{D}\cdot\frac{\partial\vec{E}}{\partial x}
\tag{32.4}
\end{equation}

\modtext{
\textbf{МОДИФИКАЦИЯ 32.3:} Тамм отбрасывает члены с $\frac{\partial}{\partial t}$ 
на этом этапе, ссылаясь на то что они войдут в $\frac{\partial\vec{g}}{\partial t}$. 
Но для бегущей волны с комплексным $k_z$ это некорректно — появляются 
дополнительные члены пропорциональные $\omega/k_z$.
}

\modblock{
\modtext{
Для динамических полей:
\begin{equation}
\modmath{
\begin{aligned}
E_x\operatorname{div}\vec{D} &= 
\frac{\partial}{\partial x}(E_x D_x) + 
\frac{\partial}{\partial y}(E_x D_y) + 
\frac{\partial}{\partial z}(E_x D_z) - 
\vec{D}\cdot\frac{\partial\vec{E}}{\partial x} + \\
&\quad + \frac{i\omega}{c}\left[\vec{D}\times\vec{H}\right]_x -
\frac{1}{c}\left[\frac{\partial\vec{D}}{\partial t}\times\vec{H}\right]_x
\end{aligned}
}
\tag{32.4$^*$}
\end{equation}
Последние два члена отсутствуют у Тамма но важны для MenDrive.
}
}

\subsection*{4. Интегрирование по объёму с разрывами}

\modtext{
\textbf{МОДИФИКАЦИЯ 32.4:} Критически важное место! Тамм применяет теорему 
Гаусса ко всему пространству сразу, предполагая что поверхностные интегралы 
на бесконечности обращаются в ноль. Но в MenDrive:
\begin{itemize}
    \item На границах $x = \pm A$ тензор $T_{ik}$ имеет скачок
    \item Частные производные $\frac{\partial T_{ik}}{\partial x_k}$ терпят разрыв
    \item Теорема Гаусса в форме (32.5) требует непрерывности частных производных
\end{itemize}
}

% \modblock{
% \modtext{
\textbf{Правильное применение теоремы Гаусса:} отдельно к каждой области:
\begin{align}
% \modmath{
\int_{V_{\text{мет}}} \operatorname{div}\mathbf{T}\,dV &= 
\oint_{S_{\text{мет}}} \mathbf{T}\cdot d\vec{S} - 
\int_{x=-A} [\mathbf{T}]\,dS \\
\int_{V_{\text{вак}}} \operatorname{div}\mathbf{T}\,dV &= 
\oint_{S_{\text{вак}}} \mathbf{T}\cdot d\vec{S} + 
\int_{x=-A} [\mathbf{T}]\,dS - 
\int_{x=+A} [\mathbf{T}]\,dS \\
\int_{V_{\text{фер}}} \operatorname{div}\mathbf{T}\,dV &= 
\oint_{S_{\text{фер}}} \mathbf{T}\cdot d\vec{S} + 
\int_{x=+A} [\mathbf{T}]\,dS
% }
\tag{32.5$^*$}
\end{align}
где $[\mathbf{T}] = \mathbf{T}^{(+)} - \mathbf{T}^{(-)}$ — скачок на границе.

\textbf{Сумма:}
\begin{equation}
\modmath{
F_x^{\text{total}} = \oint_{S_\infty} \mathbf{T}\cdot d\vec{S} + 
\underbrace{\int_{x=-A} [\mathbf{T}]\,dS + \int_{x=+A} [\mathbf{T}]\,dS}_{
\text{поверхностные члены на внутренних границах}}
}
\tag{32.6$^*$}
\end{equation}

У Тамма поверхностные члены на внутренних границах \textbf{отсутствуют} — 
это и есть критическая ошибка для MenDrive!
% }
% }

% ============================================================
% §33. ТЕНЗОР НАТЯЖЕНИЙ (стр. 169–175)
% ============================================================

\section*{\S\,33. Тензор натяжений электромагнитного поля.}

\subsection*{1. Определение тензора}

\noindent
обозначаться нами буквой $\mathbf{T}$ (без индексов). Компоненты 
тензора могут быть записаны в следующей симметричной форме:

\begin{equation}
\mathbf{T} =
\begin{pmatrix}
T_{xx} & T_{xy} & T_{xz} \\
T_{yx} & T_{yy} & T_{yz} \\
T_{zx} & T_{zy} & T_{zz}
\end{pmatrix}.
\tag{33.6}
\end{equation}

\textbf{4.} Чтобы от интегрального соотношения (33.3) между натяжениями 
и объёмными силами перейти к соотношениям дифференциальным, мы должны, 
очевидно, прежде всего преобразовать поверхностный интеграл справа 
в объёмный (или наоборот). Воспользовавшись формулами (33.2) и (33.5),

\modtext{
\textbf{МОДИФИКАЦИЯ 33.1:} Преобразование корректно только если компоненты 
тензора непрерывно дифференцируемы. В MenDrive на границах $x = \pm A$ 
это не выполняется.
}

\modblock{
\modtext{
\begin{equation}
\modmath{
\int_V \frac{\partial T_{ik}}{\partial x_k} dV = 
\oint_S T_{ik} n_k dS - \sum_{\text{внутр.}} \int_{S_j} [T_{ik}] n_k dS
}
\tag{33.7$^*$}
\end{equation}
Последний член (сумма по внутренним границам) отсутствует у Тамма.
}
}

\subsection*{2. Дифференциальные соотношения}

\noindent
получим:

\begin{equation}
f_x = \frac{\partial T_{xx}}{\partial x} + 
\frac{\partial T_{xy}}{\partial y} + 
\frac{\partial T_{xz}}{\partial z},
\quad
f_y = \frac{\partial T_{yx}}{\partial x} + 
\frac{\partial T_{yy}}{\partial y} + 
\frac{\partial T_{yz}}{\partial z},
\quad
f_z = \frac{\partial T_{zx}}{\partial x} + 
\frac{\partial T_{zy}}{\partial y} + 
\frac{\partial T_{zz}}{\partial z}.
\tag{33.7}
\end{equation}

Эти формулы и устанавливают искомые дифференциальные соотношения между 
плотностью объёмных сил $\vec{f}$ и компонентами тензора натяжений 
$\mathbf{T}$.

\textbf{5.} Из (33.7) следует, что плотность объёмных сил определяется не

\modtext{
\textbf{МОДИФИКАЦИЯ 33.2:} Для динамических полей нужно добавить член с 
импульсом поля:
}

\modblock{
\modtext{
\begin{equation}
\modmath{
f_i + \frac{\partial g_i}{\partial t} = \sum_k \frac{\partial T_{ik}}{\partial x_k}
}
\tag{33.7$^{**}$}
\end{equation}
где $\vec{g} = \frac{1}{4\pi c}[\vec{E}\times\vec{H}]$ — плотность импульса поля.
}
}

% ============================================================
% §34. ТЕНЗОР НАТЯЖЕНИЙ В ДИЭЛЕКТРИКАХ (стр. 176–185)
% ============================================================

\section*{\S\,34. Тензор натяжений в диэлектриках и магнетиках.}

\subsection*{1. Максвеллов тензор натяжений}

\noindent
Таким образом, тензор натяжений магнитного поля $\mathbf{T}'$ может быть 
получен из соответствующего тензора $\mathbf{T}'$ электрического поля 
[формула (34.2)] простой заменой $\vec{E}$ на $\vec{H}$ и 
$\varepsilon$ на $\mu$.

\modtext{
\textbf{МОДИФИКАЦИЯ 34.1:} Стандартный тензор Тамма:
}

\modblock{
\modtext{
\begin{equation}
\modmath{
T_{ik}^{\text{(Tamm)}} = \frac{1}{8\pi}\left\{
E_i D_k + E_k D_i + H_i B_k + H_k B_i - 
\delta_{ik}(\vec{E}\cdot\vec{D} + \vec{H}\cdot\vec{B})
\right\}
}
\tag{34.1$^*$}
\end{equation}
}
}

\subsection*{2. Модифицированный тензор для MenDrive}

\modtext{
\textbf{МОДИФИКАЦИЯ 34.2:} Для учёта конвективных сил на намагниченность 
и поляризацию вводится модификация \textit{только для компоненты $T_{xx}$}:
}

% \modblock{
% \modtext{
\begin{equation}
\modmath{
T_{xx}^{\text{(modified)}} = T_{xx}^{\text{(Tamm)}} + \Delta T_{xx}
}
\tag{34.2$^*$}
\end{equation}
где добавка:
\begin{equation}
\modmath{
\Delta T_{xx} = -\frac{1}{8\pi}\left[
(\vec{B}-\vec{H})\cdot\vec{H} + 
(\vec{D}-\vec{E})\cdot\vec{E}
\right]
}
\tag{34.3$^*$}
\end{equation}
а скалярные произведения:
\begin{align}
% \modmath{
(\vec{B}-\vec{H})\cdot\vec{H} &= 
(B_x - H_x)H_x + (B_y - H_y)H_y + (B_z - H_z)H_z \\
(\vec{D}-\vec{E})\cdot\vec{E} &= 
(D_x - E_x)E_x + (D_y - E_y)E_y + (D_z - E_z)E_z
% }
\end{align}

\textbf{Важные замечания:}
\begin{enumerate}
    \item Добавка $\Delta T_{xx}$ присутствует \textit{только} в компоненте 
    $T_{xx}$, так как конвективная сила действует в направлении $x$ 
    (нормаль к границам раздела сред).
    
    \item Выражения не содержат явных $\varepsilon_{ii}$ и $\mu_{ii}$ — 
    проницаемости ``спрятаны'' в соотношениях между $\vec{D}$ и $\vec{E}$, 
    $\vec{B}$ и $\vec{H}$.
    
    \item Для линейных сред:
    \begin{align}
%     \modmath{
    (\vec{B}-\vec{H})\cdot\vec{H} &= 
    (\mu_{xx}-1)H_x^2 + (\mu_{yy}-1)H_y^2 + (\mu_{zz}-1)H_z^2 \\
    (\vec{D}-\vec{E})\cdot\vec{E} &= 
    (\varepsilon_{xx}-1)E_x^2 + (\varepsilon_{yy}-1)E_y^2 + 
    (\varepsilon_{zz}-1)E_z^2
%     }
    \end{align}
\end{enumerate}
% }
% }

\subsection*{3. Матричное представление}

\modtext{
\textbf{МОДИФИКАЦИЯ 34.3:} Разложение на электрическую и магнитную части 
для наглядности:
}

% \modblock{
% \modtext{
\textbf{Электрическая часть модифицированного тензора:}
\begin{equation}
\modmath{
\mathbf{T}^{(e,\text{mod})} = \frac{1}{8\pi}
\begin{pmatrix}
2E_x D_x - (\vec{E}\cdot\vec{D}) - (\vec{D}-\vec{E})\cdot\vec{E} & 
E_x D_y + E_y D_x & 
E_x D_z + E_z D_x \\
E_y D_x + E_x D_y & 
2E_y D_y - (\vec{E}\cdot\vec{D}) & 
E_y D_z + E_z D_y \\
E_z D_x + E_x D_z & 
E_z D_y + E_y D_z & 
2E_z D_z - (\vec{E}\cdot\vec{D})
\end{pmatrix}
}
\tag{34.4$^*$}
\end{equation}

\textbf{Магнитная часть модифицированного тензора:}
\begin{equation}
\modmath{
\mathbf{T}^{(m,\text{mod})} = \frac{1}{8\pi}
\begin{pmatrix}
2H_x B_x - (\vec{H}\cdot\vec{B}) - (\vec{B}-\vec{H})\cdot\vec{H} & 
H_x B_y + H_y B_x & 
H_x B_z + H_z B_x \\
H_y B_x + H_x B_y & 
2H_y B_y - (\vec{H}\cdot\vec{B}) & 
H_y B_z + H_z B_y \\
H_z B_x + H_x B_z & 
H_z B_y + H_y B_z & 
2H_z B_z - (\vec{H}\cdot\vec{B})
\end{pmatrix}
}
\tag{34.5$^*$}
\end{equation}
% }
% }

\subsection*{4. Компонента $T_{xx}$ для TE-волны}

\modtext{
\textbf{МОДИФИКАЦИЯ 34.4:} Для TE-волны в модели MenDrive 
($E_x = E_z = H_y = 0$, ненулевые $E_y, H_x, H_z$):
}

% \modblock{
% \modtext{
\begin{equation}
\modmath{
\begin{aligned}
T_{xx}^{\text{(Tamm)}} &= \frac{1}{8\pi}\left(
2H_x B_x - E_y D_y - H_x B_x - H_z B_z
\right) \\
&= \frac{1}{8\pi}\left(
H_x B_x - E_y D_y - H_z B_z
\right)
\end{aligned}
}
\tag{34.6$^*$}
\end{equation}

Модифицированная компонента:
\begin{equation}
\modmath{
\begin{aligned}
T_{xx}^{\text{(modified)}} &= T_{xx}^{\text{(Tamm)}} 
- \frac{1}{8\pi}\left[
(B_x - H_x)H_x + (B_z - H_z)H_z + (D_y - E_y)E_y
\right] \\
&= \frac{1}{8\pi}\left(
H_x^2 - E_y^2 - H_z^2
\right)
\end{aligned}
}
\tag{34.7$^*$}
\end{equation}

\textbf{Физический смысл:} после модификации $T_{xx}^{\text{(modified)}}$ 
выражается \textit{только через напряжённости} $\vec{H}$ и $\vec{E}$, 
без явных $\vec{B}$ и $\vec{D}$. Это обеспечивает правильное воспроизведение 
конвективной силы $t_{I,\text{conv},Hn}$ при вычислении поверхностного 
интеграла через скачок на границе раздела сред.
% }
% }

% ============================================================
% §84. ТЕНЗОР НАТЯЖЕНИЙ В МАГНЕТИКАХ (стр. 386–393)
% ============================================================

\section*{\S\,84. Тензор натяжений в магнетиках.}

\subsection*{1. Магнитная часть тензора}

\noindent
движения электронов в атоме, причём это изменение по порядку величины сравнимо 
с абсолютной величиной момента атома.

\modtext{
\textbf{МОДИФИКАЦИЯ 84.1:} В модели MenDrive феррит описывается \textit{диагональным 
тензором} магнитной проницаемости:
}

\modblock{
\modtext{
\begin{equation}
\modmath{
\overleftrightarrow{\mu} = 
\begin{pmatrix}
\mu_{xx} & 0 & 0 \\
0 & \mu_{yy} & 0 \\
0 & 0 & \mu_{zz}
\end{pmatrix}
}
\tag{84.1$^*$}
\end{equation}
Это означает, что соотношения между полями имеют вид:
\begin{equation}
\modmath{
B_i = \mu_{ii} H_i, \quad D_i = \varepsilon_{ii} E_i \quad (\text{без суммирования по } i)
}
\tag{84.2$^*$}
\end{equation}
}
}

\subsection*{2. Плотность силы для тензорных сред}

\modtext{
\textbf{МОДИФИКАЦИЯ 84.2:} Для диагональных тензоров проницаемостей конвективная 
сила на намагниченность принимает вид:
}

% \modblock{
% \modtext{
\begin{equation}
\modmath{
\left[(\vec{M}\cdot\nabla)\vec{H}\right]_x = 
\frac{\mu_{xx} - 1}{4\pi}H_x\frac{\partial H_x}{\partial x} + 
\frac{\mu_{yy} - 1}{4\pi}H_y\frac{\partial H_y}{\partial x} + 
\frac{\mu_{zz} - 1}{4\pi}H_z\frac{\partial H_z}{\partial x}
}
\tag{84.3$^*$}
\end{equation}

Интегрируя по глубине материала и учитывая затухание поля на бесконечности:
\begin{equation}
\modmath{
F_x^{(\text{conv})} = 
-\frac{\mu_{xx} - 1}{8\pi}|H_x(A)|^2 
-\frac{\mu_{yy} - 1}{8\pi}|H_y(A)|^2 
-\frac{\mu_{zz} - 1}{8\pi}|H_z(A)|^2
}
\tag{84.4$^*$}
\end{equation}

\textbf{Важно:} для TE-волны в вашей модели ($H_y = 0$, ненулевые $H_x, H_z$):
\begin{equation}
\modmath{
F_x^{(\text{conv, TE})} = 
-\frac{\mu_{xx} - 1}{8\pi}|H_x(A)|^2 
-\frac{\mu_{zz} - 1}{8\pi}|H_z(A)|^2
}
\tag{84.5$^*$}
\end{equation}
% }
% }

% ============================================================
% §105. ТЕНЗОР НАТЯЖЕНИЙ И ПОНДЕРОМОТОРНЫЕ СИЛЫ (стр. 506–511)
% ============================================================

\section*{\S\,105. Тензор натяжений и пондеромоторные силы электромагнитного поля.}

\subsection*{1. Обобщение на случай переменного поля}

\noindent
\textbf{1.} В §§\,34 и 84 мы нашли выражение тензора натяжений 
$\mathbf{T}$ в стационарных электрическом и магнитном полях. Покажем 
теперь, что если допустить применимость этих выражений и в случае 
произвольного переменного электромагнитного поля, то вытекающие из 
этого допущения выводы будут находиться в полном соответствии с 
результатами настоящей главы.

\modtext{
\textbf{МОДИФИКАЦИЯ 105.1:} Тамм \textit{постулирует} применимость статического 
тензора к динамике. Но это требует доказательства через явный вывод с учётом 
$\frac{\partial}{\partial t}$.
}

\modblock{
\modtext{
Для динамических полей тензор натяжений должен включать дополнительные члены:
\begin{equation}
\modmath{
\mathbf{T}^{\text{(dynamic)}} = \mathbf{T}^{\text{(static)}} + 
\Delta\mathbf{T}^{\text{(dynamic)}} + 
\Delta\mathbf{T}^{\text{(magnetic)}} + 
\Delta\mathbf{T}^{\text{(gradient)}}
}
\tag{105.1$^*$}
\end{equation}
}
}

\subsection*{2. Уравнение баланса импульса}

\noindent
В §§\,34 и 84 мы убедились, что тензор напряжений электромагнитного 
поля $\mathbf{T}$ может быть разложен на сумму «максвеллова» тензора 
$\mathbf{T}'$ и тензора стрикционных натяжений $\mathbf{T}''$. В 
дальнейшем мы для простоты вовсе не будем рассматривать стрикционных 
натяжений $\mathbf{T}''$, приводящих лишь к перераспределению 
пондеромоторных сил по объёму находящихся в поле тел, но не изменяющих 
результирующей этих сил, и будем отождествлять полный тензор натяжений 
$\mathbf{T}$ с максвелловым тензором $\mathbf{T}'$. При этих

\modtext{
\textbf{МОДИФИКАЦИЯ 105.2:} Тамм выводит уравнение (105.9):
}

\modblock{
\modtext{
\begin{equation}
\modmath{
f_i + \frac{\partial g_i}{\partial t} = \sum_k \frac{\partial T_{ik}}{\partial x_k}
}
\tag{105.9$^*$}
\end{equation}
где $\vec{g} = \frac{1}{4\pi c}[\vec{E}\times\vec{H}]$ — плотность электромагнитного 
импульса.
}
}

\subsection*{3. Интегрирование по всему пространству}

\noindent
движения $\vec{g}$. Этот 
дополнительный член обеспечивает сохранение суммы количества движения 
материальных тел и электромагнитного поля.

Действительно, проинтегрируем уравнение (105.9) по всему пространству, 
предполагая, что слагающие электромагнитного тензора натяжений достаточно 
быстро убывают по мере удаления в бесконечность. Интеграл правой части 
уравнения (105.9) при этих усло-

\noindent
виях обратится в нуль, ибо, например, при $i = 1$

\begin{equation}
\sum_k\frac{\partial T_{xk}}{\partial x_k}
\end{equation}

\modtext{
\textbf{МОДИФИКАЦИЯ 105.3:} Критическое место! Тамм явно пишет условие 
применимости своего вывода на стр.~510--511:
}

% \modblock{
% \modtext{
\begin{quote}
\modtext{
\textit{«Действительно, проинтегрируем уравнение (105.9) по всему пространству, 
предполагая, что слагающие электромагнитного тензора натяжений достаточно 
быстро убывают по мере удаления в бесконечность. Интеграл правой части 
уравнения (105.9) при этих условиях обратится в нуль»}
}
\end{quote}

И далее:
\begin{quote}
\modtext{
\textit{«где поверхностный интеграл должен быть взят по поверхности $S$ объёма $V$ 
и при удалении этой поверхности в бесконечность обратиться, по предположению, в нуль»}
}
\end{quote}

\textbf{В модели MenDrive это условие НЕ выполняется:}
\begin{enumerate}
    \item На внутренних границах $x = \pm A$ тензор $T_{ik}$ имеет скачок
    \item При $k_y = 0$ поле не убывает вдоль $y$ (бесконечный резонатор)
    \item Для бегущей волны с комплексным $k_z$ условие $T_{ik} \to 0$ при 
    $z \to \infty$ требует проверки
\end{enumerate}
% }
% }

\subsection*{4. Закон сохранения импульса}

\noindent
Отсюда закон сохранения (105.10):

\begin{equation}
\frac{d}{dt}\left(\vec{G}_{\text{мех}} + \int \vec{g}\,dV\right) = 0
\tag{105.10}
\end{equation}

\modtext{
\textbf{МОДИФИКАЦИЯ 105.4:} Для MenDrive с внутренними границами:
}

\modblock{
\modtext{
\begin{equation}
\modmath{
\frac{d}{dt}\left(\vec{P}_{\text{мех}} + \vec{P}_{\text{Пойнтинг}} + 
\vec{P}_{\text{вакуум}}\right) = 0
}
\tag{105.10$^*$}
\end{equation}
где:
\begin{equation}
\modmath{
\vec{P}_{\text{вакуум}} = \frac{1}{c} \sum_{j} \int_{S_j} \mathbf{T} \cdot d\vec{S}_j
}
\tag{105.11$^*$}
\end{equation}
— «импульс вакуума» от внутренних поверхностных интегралов на границах 
раздела сред.
}
}

\subsection*{5. Баланс для MenDrive}

\modtext{
\textbf{МОДИФИКАЦИЯ 105.5:} Полный баланс с учётом всех областей:
}

% \modblock{
% \modtext{
\begin{equation}
\modmath{
F_x^{\text{total}} = F_x^{\text{мет}} + F_x^{\text{вак}} + F_x^{\text{фер}} + F_x^{\text{wire}}
}
\tag{105.12$^*$}
\end{equation}
где:
\begin{align}
% \modmath{
F_x^{\text{мет}} &= -T_{xx}^{\text{(modified)}}\big|_{x=-A^-} \\
F_x^{\text{вак}} &= T_{xx}^{\text{(modified)}}\big|_{x=-A^+} - T_{xx}^{\text{(modified)}}\big|_{x=+A^-} \\
F_x^{\text{фер}} &= +T_{xx}^{\text{(modified)}}\big|_{x=+A^+} \\
F_x^{\text{wire}} &= \text{сила на виток возбуждения}
% }
\end{align}

Для стационарного режима:
\begin{equation}
\modmath{
\langle F_x^{\text{total}} \rangle = 0
}
\tag{105.13$^*$}
\end{equation}
что соответствует закону сохранения импульса, но отдельные компоненты 
$\langle F_x^{\text{мет}} \rangle$, $\langle F_x^{\text{фер}} \rangle$ могут быть 
сонаправлены, создавая результирующую тягу на резонатор.
% }
% }

% ============================================================
% ЗАКЛЮЧЕНИЕ
% ============================================================

\section*{Заключение}

% \modtext{
Приведённые модификации позволяют:
\begin{enumerate}
    \item Корректно учесть динамические эффекты в MenDrive (формула 18.9 
    вместо 17.5)
    \item Включить вклад магнитных токов в феррите ($\vec{j}_m \neq 0$)
    \item Правильно обработать разрывы $\varepsilon$, $\mu$ на границах 
    (теорема Гаусса отдельно к каждой области)
    \item Получить модифицированный тензор, воспроизводящий конвективные 
    силы $t_{I,\text{conv},Hn}$
    \item Согласовать поверхностный интеграл с объёмным расчётом сил
\end{enumerate}

\textbf{Важно:} Все модификации выделены \textcolor{modcolor}{зелёным цветом}. 
Для изменения цвета достаточно поменять определение \texttt{modcolor} в преамбуле.

\textbf{Физический итог:} Стандартный тензор Тамма неполон для симметричной 
системы уравнений Максвелла с магнитными токами. Модификация через добавку 
$\Delta T_{xx}$ позволяет воспроизвести конвективную силу на намагниченность 
через поверхностный интеграл, что согласует два подхода к вычислению силы.
% }

\end{document}
